% QUESTAO
Indicamos genericamente uma matriz $A$ de ordem $m\times n$, ou seja, como $m$ linhas e $n$ colunas, por $A=(a_{ij})_{m\times n}$. No elemento $a_{ij}$, o índice $i$ indica a linha, e o $j$, a coluna em que o elemento está localizado. O elemento $a_{32}$, por exemplo, tem $i=3$ e $j=2$, ou seja, está localizado na 3ª linha e na 2ª coluna. Assim, supondo uma matriz $A=(a_{ij})_{3\times 5}$, tal que $A = \left[\begin{array}{ccccc}
0 & -1 & 5{,}7 & -3 & 2 \\
4 & 15 & 3 & \sqrt{7} & 6 \\
\pi & -9 & 10 & -2 & 8
\end{array}\right]$. Qual o valor dos elementos $a_{11}$, $a_{13}$, $a_{31}$ e $a_{24}$, respectivamente?

% ALTERNATIVAS
$0$; $5{,}7$; $\pi$ e $\sqrt{7}$
$0$; $\pi$; $5{,}7$ e $\sqrt{7}$
$0$; $5{,}7$; $\pi$ e $10$
$8$; $5{,}7$; $15$ e $\sqrt{7}$
$\pi$; $5{,}7$; $\pi$ e $-2$


